% HANFIS - Results: observability overhead (auto-generated numbers from
% results_python/telemetry/paired_effects_v3.csv, 2026-09-22).
% Ten counterbalanced pairs, protocol external-v3.0, Runs 7-16.
% Exact two-sided sign-flip test (floor p = 2/2^10 = 0.002), 10,000-resample bootstrap CI.

\subsection{Observability overhead}
\label{subsec:overhead}

Ten complete run pairs passed every acceptance gate (Runs 7--16, protocol
\texttt{external-v3.0}, counterbalanced five ON--OFF and five OFF--ON). One aborted
attempt of Run~7 and one aborted attempt of Run~17 failed the integrity gate and were
recorded as incidents rather than observations, in accordance with the fail-closed
rule stated in Section~\ref{sec:methodology}. Table~\ref{tab:overhead} reports the
paired differences $\Delta Y_r = \bar{Y}_{r,\mathrm{ON}} - \bar{Y}_{r,\mathrm{OFF}}$.

\begin{table}[htbp]
\caption{Paired monitoring overhead over ten counterbalanced runs (protocol external-v3.0).
CI: 10{,}000-resample bootstrap; $p$: exact two-sided sign-flip test (floor $2/2^{10}=0.002$).}
\label{tab:overhead}
\begin{tabular}{@{}lrrrr@{}}
\toprule
Measure & Mean $\Delta$ & 95\% CI & Relative & $p$ \\
\midrule
\multicolumn{5}{@{}l}{\emph{Resource cost of the observability stack}}\\
Node CPU (mCPU)            & $+99.36$  & $[96.29,\ 103.35]$ & $+43.3\%$ & $0.002$ \\
Monitoring namespace CPU (mCPU) & $+39.79$ & $[33.00,\ 52.77]$ & --- & $0.002$ \\
Node memory (MiB)          & $+280.76$ & $[250.99,\ 303.54]$ & $+18.7\%$ & $0.002$ \\
Monitoring namespace memory (MiB) & $+204.19$ & $[191.38,\ 218.82]$ & --- & $0.002$ \\
\multicolumn{5}{@{}l}{\emph{Effect on the observed workload}}\\
Latency median (ms)        & $-0.005$  & $[-0.055,\ +0.044]$ & $-0.06\%$ & $0.859$ \\
Latency P95 (ms)           & $+0.059$  & $[-0.126,\ +0.265]$ & $+0.67\%$ & $0.607$ \\
Latency P99 (ms)           & $+0.188$  & $[-0.127,\ +0.448]$ & $+1.94\%$ & $0.268$ \\
P0 container CPU (mCPU)    & $-0.059$  & $[-0.156,\ +0.041]$ & $-3.9\%$  & $0.305$ \\
P0 container memory (MiB)  & $-0.334$  & $[-0.522,\ -0.147]$ & $-1.19\%$ & $0.016$ \\
\botrule
\end{tabular}
\end{table}

Two results follow. First, the resource cost of the observability stack is large,
consistent, and unambiguous: every pair shows the same direction and the effect sizes
are extreme ($d_z = 16.5$ for node CPU). Node-level CPU increases by $99.4$~mCPU while
the monitoring namespace itself accounts for only $39.8$~mCPU, so pod-level accounting
understates the true cost by a factor of approximately $2.5$; the remainder is induced
platform work (kubelet, cAdvisor, container runtime, and the API server serving the
additional scrape traffic). The same pattern holds for memory ($280.8$~MiB at node
level against $204.2$~MiB inside the namespace).

Second, this resource cost does not propagate to the observed workload. For median
latency and P95 latency we fail to reject the null hypothesis, with confidence
intervals of $[-0.06,+0.04]$~ms and $[-0.13,+0.27]$~ms respectively, that is, below
$3\%$ of the baseline latency. The P0 container consumes marginally \emph{less}
memory under monitoring ($-0.33$~MiB, $p = 0.016$), an effect of practical size zero
that we report for completeness rather than interpret. With ten pairs the sign-flip
test can reach $p = 0.002$, so the failure to reject on latency reflects the absence
of a detectable effect rather than the absence of statistical power.

These measurements justify treating the observability stack as a fixed resource
overhead that must be budgeted for, and justify using the collected telemetry as
HANFIS features without correcting the workload latency for instrumentation effects.
We do not claim equivalence: the confidence intervals bound the effect, they do not
exclude it.
